% !TEX root = ../main.tex
% ============================================================
% RESUMEN: se escribe al final, cuando los resultados estén.
% ============================================================
\chapter*{Resumen}

% Repo: los números, del capítulo 5 (y por él, del 4).
Reconstruir un holograma en microscopía holográfica digital sin lentes (DLHM)
exige propagar numéricamente el campo entre el sensor y la muestra. El
espectro angular calculado con la FFT, el método más usado, impone a todo el
cálculo el muestreo de la entrada; el espectro angular por producto matricial
(MPASM) lo libera con dos coeficientes de control, $K_f$ y $s$, pero solo se
había evaluado en simulación. En este trabajo se implementaron el FFT-ASM, el
método de banda limitada y el MPASM en un mismo paquete de Python para CPU y
GPU, verificado contra una solución analítica y contra la integral de
Rayleigh--Sommerfeld; se formuló un MPASM que lleva dentro de sus matrices la
iluminación esférica exacta de la DLHM; se midió la validez de los métodos en
la geometría del banco del laboratorio y en una ida y vuelta desde distancias
cortas hasta la zona de Fresnel, y se reconstruyeron hologramas experimentales
de \emph{targets} USAF de amplitud y de fase, cuya geometría se infirió de los
propios hologramas. El control del muestreo rinde cuando la distancia de
propagación equivalente es del orden de $N(\delta/M)^2/\lambda$ o mayor, es
decir, cuando el campo no cabe en la ventana del registro, lo que en la DLHM
ocurre con la muestra cerca de la fuente y no a distancias físicas grandes. En
ese régimen, el MPASM con sobremuestreo propaga con entre 11 y 1100 veces
menos error que el FFT-ASM con relleno. Al reconstruir, con la misma ventana,
el MPASM esférico y el método del laboratorio dan casi la misma imagen
(correlación de 0.87 a 0.97). La ventaja del MPASM es de memoria: en una GPU
de \SI{4}{GB} procesa un recorte del sensor veinte veces mayor, con el que el
contraste de los detalles de 5 a \SI{10}{\micro\metre} casi se duplica.

\vspace{0.5cm}
\noindent\textbf{Palabras clave:} microscopía holográfica digital sin lentes;
espectro angular; producto matricial; muestreo; propagación numérica;
reconstrucción de hologramas.

% Abstract opcional: las directrices solo piden el resumen.
% \chapter*{Abstract}
% ...
% \noindent\textbf{Keywords:} digital lensless holographic microscopy;
% angular spectrum; matrix product; sampling; numerical propagation.
